\section{Transfer Accounting and Policy Counterfactuals}
\label{sec:counterfactuals}

The estimates support a compact transfer accounting, and its striking
feature is how little of it is assumption. Table~\ref{tab:welfare}
collects the rows.

\begin{table}[!htbp]
  \centering
  \singlespacing
  \caption{\capttitle{The Tip-Tax Ledger (USD Millions per Year)}}
  \label{tab:welfare}
  \small
  \begin{tabular}{@{}p{0.72\textwidth}r@{}}
    \toprule
    & USD m/yr \\
    \midrule
    \input output/tables/tab-welfare.tex
    \bottomrule
  \end{tabular}
  \floatnotes{Transfers are sums of observed fee payments on
  standard-rate card trips (2025 toll: 2025H1 annualized; 2019
  surcharge: calendar 2019) multiplied by the design-estimated
  pass-through. The counterfactual rows apply the dose linearity
  established across the \$0.50--\$2.50 range. The deadweight-loss bound
  is one-half times the toll times the trip decline at the lower edge of
  the demand confidence interval; the point estimate of that decline is
  indistinguishable from zero.}
\end{table}

\subsection{Transfers}
The toll moves \$16.6 million a year from riders to the MTA on card
trips---the intended, legislated flow---and \$2.2 million a year from
riders to drivers through the tip base, the unintended flow this paper
measures. The 2019 congestion surcharge, at its larger dose and volume,
moved about \$18.7 million a year through the same channel. These are
transfers, not losses: dollars change hands at face value.

\subsection{Counterfactual designs}
Because pass-through is stable in the dose from \$0.50 to \$5.00---the
JFK flat-fare surcharge supplies the \$4.50 and \$5.00 points
(Section~\ref{sec:results-controls})---the formula
$\rho \times \text{fee} \times \text{volume}$ prices design
alternatives directly. A toll set at \$2.50 would have induced roughly
\$7.4 million a year in tips. A mandated fare-only tip base---the
design one payment vendor used voluntarily in 2012
(Appendix~\ref{app:replication})---sets the induced flow to zero at no
collection cost, returning the full \$21 million a year across both
fees to riders. The tip-side of any future meter fee can be priced the
same way before enactment.

\subsection{The base-rule menu}
Because pass-through applies to whatever the base contains, every
candidate base rule has a price, and Table~\ref{tab:menu} prices the
2025 fee environment. On card trips, \$210 million a year of non-fare
dollars sit inside the tip base, generating \$23--32 million a year in
tips at the paper's two scalings. The government-remitted components
alone contribute \$141 million of base and \$15--21 million of tips---
and the largest of them is not the new CBD toll but the 2019 New York
State congestion surcharge, still running at \$2.50 per trip: \$73
million of base, \$8--11 million a year of induced tips. A regulator
choosing the base is choosing a row of this menu. Excluding the CBD
toll alone returns \$2--3 million a year to riders; excluding every
government fee returns \$15--21 million; the fare-only base that one
vendor used in 2012 returns \$23--32 million, at zero collection cost.
Tips induced by the driver-kept surcharges are a different object---a
gratuity on a price rather than on a tax---and the table labels them
separately.

\begin{table}[!htbp]
  \centering
  \singlespacing
  \caption{\capttitle{The Base-Rule Menu (USD Millions per Year, 2025
  Fee Environment)}}
  \label{tab:menu}
  \small
  \begin{tabular}{@{}p{0.52\textwidth}rcc@{}}
    \toprule
    & In tip base & \multicolumn{2}{c}{Induced tips at $\rho$ of:} \\
    \cmidrule(lr){3-4}
    Component / base rule & (USD m/yr) & $0.108$ & $0.150$ \\
    \midrule
    \input output/tables/tab-menu.tex
    \bottomrule
  \end{tabular}
  \floatnotes{Card trips, all rate codes, 2025H1 annualized. ``In tip
  base'' sums each component's dollars on card trips. Induced tips
  apply the pooled ITT per statutory dollar ($0.108$) and $\rho^{IV}$
  per recorded dollar ($0.150$); collections are recorded dollars, so
  the higher column is the aggregation-consistent one and the lower is
  conservative. Base rules give the tips returned to riders if the
  percentage buttons excluded the listed components.}
\end{table}

\subsection{Efficiency: the margins we measured are silent}
A transfer becomes a welfare problem when it distorts behavior, and on
every margin we can measure, it does not. Riders do not retime hails
around the 4:00pm surcharge (density ratio $0.99$), do not reduce
congestion-zone ridership detectably ($-1.0$ percent, SE $0.9$), do
not flee to cash (the card share rises at the cutoff), and barely
adjust tipping behavior itself (exact-default use falls under one
point). As a gauge of scale rather than a welfare bound: a Harberger
triangle computed at the unfavorable edge of the demand confidence
interval---one-half times the toll times a 2.9 percent trip
decline---comes to \$0.34 million a year, an order of magnitude below
the transfer, and the point estimate of the decline is zero. Margins we
cannot measure---whether inattentive riders would endorse the tip on
reflection, and any driver labor-supply response---are discussed below.

\subsection{What kind of transfer this is}
Two clarifications about its nature. First, on the receiving side, a
tip is the most driver-kept dollar in the fare stack: unlike metered
revenue it does not offset lease payments or commissions, so an induced
tip dollar is close to marginal driver income (net of card-processing
fees, and taxable, since card tips are recorded). The driver remits the
congestion fee itself and nets nothing on it; the induced tip is the
only part of the flow that stays. Second, the transfer is not
regulatory capture in the classic sense of an interested industry
bending rules in its favor: no beneficiary chose the base. The TLC
certifies the payment systems, so the base rule sits within the
regulable surface, but the record shows convergence by payment vendors,
not advocacy by drivers---closer to regulatory neglect than capture.
``Unlegislated'' describes a decision nobody made.

\subsection{What this accounting cannot say}
Whether the transfer harms riders depends on whether the inattentive
four-fifths would endorse the tip on reflection---a preferences
question our design-based evidence cannot answer, so we report
transfers as transfers. Driver labor supply, medallion values, and
general-equilibrium effects are likewise out of scope. The
distributional geography of the burden appears in
Section~\ref{sec:results-fee}; we label it neighborhood exposure, not
individual incidence.
